\documentclass[11pt,letterpaper,reqno]{amsart}
\usepackage{amssymb}
\usepackage{amsmath}
\usepackage{amsthm}
\usepackage{amsfonts}
\usepackage{enumitem}
\usepackage{booktabs}
\usepackage{graphicx}
\usepackage{xcolor}
\usepackage[T1]{fontenc}
\usepackage{doi}
\usepackage{comment}
\usepackage{pgfplots}
\pgfplotsset{compat=1.18}
\addtolength{\hoffset}{-1.5cm}\addtolength{\textwidth}{3cm}
\addtolength{\voffset}{-1cm}\addtolength{\textheight}{2cm}
\usepackage{hyperref}
\hypersetup{pdfstartview={FitH}}
\usepackage{bookmark}
\usepackage[capitalize,noabbrev]{cleveref}

\newtheorem{theorem}{Theorem}[section]
\newtheorem{lemma}[theorem]{Lemma}
\newtheorem{proposition}[theorem]{Proposition}
\newtheorem{corollary}[theorem]{Corollary}
\newtheorem{claim}{Claim}
\newtheorem{question}[theorem]{Question}
\newtheorem{problem}[theorem]{Problem}
\newtheorem{conjecture}[theorem]{Conjecture}
\theoremstyle{definition}
\newtheorem{example}[theorem]{Example}
\newtheorem{remark}[theorem]{Remark}
\newtheorem{definition}[theorem]{Definition}

\newtheorem{theoremletter}{Theorem}
\renewcommand{\thetheoremletter}{\Alph{theoremletter}}
\crefname{theoremletter}{Theorem}{Theorems}
\Crefname{theoremletter}{Theorem}{Theorems}
\newcommand{\dd}{\,\mathrm d}
\DeclareMathOperator{\cof}{cof}

\begin{document}

\title{Proof of the Planar P\'olya--Szeg\H{o} Buckling Conjecture}


\author[Q.~Tang]{Quanyu Tang}
\address{School of Mathematical Sciences, University of Science and Technology of China, Hefei 230026, P.~R.~China}
\email{tangquanyu827@gmail.com}

\author[H.~Zhang]{Haiqi Zhang}
\address{School of Mathematics, Shandong University, Jinan 250100, P.~R.~China}
\email{ZHQAQ2024@outlook.com}

\date{Working draft v5 -- 18 October 2026}

\maketitle

\begin{center}
\small\textbf{Work in progress.} This manuscript is under revision
and has not yet been submitted to arXiv.
\end{center}

\section{Introduction}
For an open set $\Omega\subset\mathbb R^2$ of finite positive area, define
\[
 \Lambda_1(\Omega)=\inf_{0\ne v\in H_0^2(\Omega)}
       \frac{\int_\Omega(\Delta v)^2\dd x}
            {\int_\Omega|\nabla v|^2\dd x},\qquad
 \lambda_1^{\mathrm D}(\Omega)=\inf_{0\ne v\in H_0^1(\Omega)}
       \frac{\int_\Omega|\nabla v|^2\dd x}
            {\int_\Omega v^2\dd x}.
\]
Here $H_0^m(\Omega)$ denotes the closure of $C_c^\infty(\Omega)$
in $H^m(\Omega)$. This definition is used throughout, including when
$\partial\Omega$ is irregular. For a smooth bounded domain, the first
quantity is the least eigenvalue of the clamped buckling problem
\[
 \Delta^2v+\Lambda\Delta v=0\quad\hbox{in }\Omega,
 \qquad v=\partial_\nu v=0\quad\hbox{on }\partial\Omega,
\]
where $\nu$ is the outward unit normal. In the planar plate model,
$\Lambda_1$ represents the critical load for uniform lateral compression.

The isoperimetric question for this eigenvalue belongs to the classical
program of comparing membranes and plates of prescribed area. The
Faber--Krahn theorem \cite{Faber,Krahn} identifies the disk as the
minimizer of $\lambda_1^{\mathrm D}$. P\'olya and Szeg\H{o}
\cite{PS} proposed the corresponding assertion for the first buckling
load. Let $J_\nu$ be the Bessel function of the first kind, and let
$j_{\nu,n}$ denote its $n$th positive zero. On a disk of radius $R$,
\[
 \lambda_1^{\mathrm D}(B_R)=j_{0,1}^2/R^2,
 \qquad \Lambda_1(B_R)=j_{1,1}^2/R^2.
\]
The planar P\'olya--Szeg\H{o} conjecture takes the following form.
\begin{conjecture}\label{conj:PS}
For every smooth bounded planar domain $\Omega$,
\[
 |\Omega|\Lambda_1(\Omega)\ge\pi j_{1,1}^2,
\]
with equality if and only if $\Omega$ is a disk.
\end{conjecture}

Szeg\H{o}'s early symmetrization argument \cite{Szego50,Szego58}
established the buckling inequality under an additional hypothesis
that a first eigenfunction does not change sign. The correction in
\cite{Szego58} addresses the admissibility of the rearranged function.
The sign hypothesis cannot be imposed on general buckling eigenfunctions;
see the discussion in \cite{AL} and the analysis of annuli in \cite{BP}.
Thus the usual positivity-based proof for the Dirichlet Laplacian does
not directly settle Conjecture~\ref{conj:PS}.

A second line of work compares buckling and membrane spectra. Payne's
inequality $\Lambda_1\ge\lambda_2^{\mathrm D}$ \cite{Payne}, with the
correction discussed by Friedlander \cite{Friedlander}, gives a useful
but nonsharp area bound. Here $\lambda_j^{\mathrm D}$ denotes the
$j$th Dirichlet eigenvalue, counted with multiplicity. In dimension two, the bound of Bramble and
Payne \cite{BramblePayne} is $|\Omega|\Lambda_1>2\pi j_{0,1}^2$ for
bounded connected domains. Ashbaugh and Laugesen \cite{AL} extended this
estimate to higher dimensions. More recently, Levitin and Polterovich
\cite{LP} established Payne's shifted comparison for every odd buckling
index. These index comparisons are distinct from the comparison with
$\lambda_1^{\mathrm D}$ proved below.

Shape optimization provides another approach. Ashbaugh and Bucur
\cite{AB} proved existence results in the simply connected planar class
and for a relaxed problem. Stollenwerk \cite{Stollenwerk} subsequently
proved existence among finite-measure open sets in arbitrary dimension.
The Weinberger--Willms argument, presented in \cite{AB}, identifies a
smooth simply connected optimizer as a disk; related optimality
conditions were developed by Stollenwerk and Wagner \cite{SW}.
Lamboley and Nahon \cite{LN} proved boundedness and partial boundary
regularity of planar optimizers, and weakened the regularity assumptions
in this characterization. Those results leave hypotheses on the
optimizer; the argument here does not require an optimal-domain
regularity theorem.

Our main result is a sharp comparison with the first Dirichlet eigenvalue.
\begin{theorem}\label{main}
For every open set $\Omega\subset\mathbb R^2$ with
$0<|\Omega|<\infty$,
\begin{equation}\label{eq:comparison}
 \Lambda_1(\Omega)\ge
       \frac{j_{1,1}^2}{j_{0,1}^2}\lambda_1^{\mathrm D}(\Omega).
\end{equation}
The constant is optimal.
\end{theorem}

Combining Theorem~\ref{main} with the Faber--Krahn inequality resolves
Conjecture~\ref{conj:PS}, without assumptions on topology or boundary
regularity.
\begin{corollary}\label{cor:area}
For every open set $\Omega\subset\mathbb R^2$ of finite positive area,
\begin{equation}\label{eq:area}
 |\Omega|\Lambda_1(\Omega)\ge\pi j_{1,1}^2.
\end{equation}
Equality holds if and only if $\Omega$ is an open disk as a set.
\end{corollary}

The proof of \eqref{eq:comparison} is variational. We construct a
one-variable function $f$ for which
\[
 F=|\nabla U|f\bigl(U/|\nabla U|\bigr)
\]
converts a compactly supported buckling test function into a Dirichlet
test function. The value at $\nabla U=0$ is specified by continuity.
After adding two divergence terms whose integrals vanish, the desired
functional inequality reduces to positivity of a symmetric matrix.
Its kernel is determined by the radial disk eigenfunctions. The
remaining $2\times2$ block is controlled by convexity, two Wronskian
comparisons, and a residue expansion for an auxiliary Bessel function.
The construction is compatible with changes of sign of $U$. In
particular, it does not use positivity of a buckling eigenfunction.

Section~\ref{sec:calibration} defines the profiles and proves the matrix
identities. Sections~\ref{cv:sec:convexity}--\ref{sec:G} establish the
one-variable inequalities. Section~\ref{sec:functional} treats the weak
interpretation of the divergence terms and proves Theorem~\ref{main}.
Section~\ref{sec:equality} proves the equality statement. The
appendices collect the fixed Bessel estimates and the form-domain facts
used in the proof. All functions may be taken real-valued: the quotient
of a complex function is a weighted average of the quotients of its
nonzero real and imaginary parts.



\subsection*{Lean formalization}

A Lean 4 formalization of the main results is available at
\url{https://github.com/QuanyuTang/polya-szego-planar-unconditional-palomar}.
The formal development proves the sharp buckling--Dirichlet comparison
in Theorem~\ref{main}, the area inequality in Corollary~\ref{cor:area},
the exact characterization of equality by an open disk, and the
buckling eigenvalue formula for disks. Its variational definitions
use the closures of smooth compactly supported functions in the
corresponding Sobolev norms, including for irregular open sets.

The formalization has been verified and registered in the
\href{https://palomar-registry.org/entry?id=PALOMAR-2026-10-18-000008&version=1}
{Palomar Registry} (entry PALOMAR-2026-10-18-000008, version~1).
The repository provides the complete Lean sources, pinned dependencies,
reproducible build instructions, verification records, and an audit
of the axioms used by the final declarations.

The formalization was developed from the authors' original informal
proof. It retains the central functional comparison based on Bessel
profiles, divergence identities, and matrix positivity, but it is
not a line-by-line transcription of the argument in this paper.
In particular, the formal development includes separate proofs of
the Faber--Krahn inequality, its equality characterization, and the
additional rigidity needed for exact equality of open sets.
The formal verification establishes the stated mathematical results,
rather than verifying every individual step of the informal
proof presented here.


\subsection*{Use of artificial intelligence}

The mathematical strategy and the original informal proof presented
in this paper were developed by the authors with assistance from ChatGPT in exploring technical arguments and preparing the
manuscript. Harmonic's Aristotle and OpenAI's Codex constructed the Lean formalization
based on the authors' original informal proof. The formal development
was subsequently refined and checked using Lean 4; as explained
above, it is not a line-by-line formalization of this manuscript.
The authors take responsibility for the mathematical statements,
the informal proof, and the accuracy of the description of the
formalization.


\section{The auxiliary profiles and matrix identity}\label{sec:calibration}
\subsection{The disk profiles}
Put $k=j_{1,1}$, $q=j_{0,1}$, $\alpha=q/k$,
$\lambda=\alpha^2$, $c=-J_0(k)>0$, and, for $0<r<k$,
\[
 \begin{gathered}
 p=J_1(r),\qquad u=J_0(r)+c,\qquad \phi=J_0(\alpha r),\\
 A=-p_r/p,\quad C=-1/r,\quad t=u/p,\quad T=t_r=tA-1.
 \end{gathered}
\]
Here $p,u,\phi>0$. The radial functions $u$ and $\phi$ are the
first buckling and Dirichlet eigenfunctions on $B_k$, with eigenvalues
$1$ and $\lambda$, respectively. Since $u_r=-p$, the change of variable
$t=u/p$ is chosen so that the transform $p f(t)$ sends $u$ to $\phi$.
Define
\[
 f(t(r))=\phi/p,\qquad z=f_t,\qquad g=f-tz,\qquad
 w=\phi_r/p,
 \qquad
 K=\left(\frac{\alpha J_1(q)}c\right)^2.
\]
A profile written as a function of $r$ means its composition with
$t=t(r)$ when necessary; derivatives carry a subscript $r$ or $t$. Primes in radial calculations
denote $r$ derivatives. All terminating decimals below denote exact
rational numbers.
The Bessel equation, defining series and recurrences are given in
\cite[Eqs.~10.2.1--10.2.2 and 10.6.1--10.6.3]{DLMF};
the canonical product and interlacing of real zeros are given in
\cite[Eqs.~10.21.15 and 10.21.2]{DLMF}.
Appendix~\ref{app:constants} proves all fixed constant bounds used below,
including the identification of their root indices.

\subsection{The change of variable}
\begin{lemma}
$T<0$ on $(0,k)$, and $t$ maps that interval bijectively onto $(0,\infty)$
in reverse order.
\end{lemma}
\begin{proof}
Let $S=p^2+up_r$, so $T=-S/p^2$. If $p_r\ge0$ then $S>0$.
The elementary alternating Bessel bounds
$J_0(r)\ge1-r^2/4$ and $J_1(r)/r\le1/2$ imply
$p_r=J_0-p/r>0$ on $0<r\le\sqrt5/2$.
If $p_r(r)<0$, necessarily $r>\sqrt5/2$. For $L=p_r/p$, the Bessel
ODE gives
\[
 (\log p)_{rr}
 =-1+\frac1{r^2}-\frac Lr-L^2
 =-1+\frac5{4r^2}-\left(L+\frac1{2r}\right)^2<0
 \quad(r>\sqrt5/2).
\]
Thus $\log p$ is strictly concave on $[r,k)$ and, since
$A(r)>0$,
\[
 u(r)=\int_r^k p(s)\,ds
 <\int_r^\infty p(r)e^{-A(r)(s-r)}ds=\frac{p(r)}{A(r)}.
\]
It follows that $S=p^2(1-tA)>0$. The endpoint expansions
$t\sim2(1+c)/r$ at zero and $t\sim(k-r)/2$ at $k$ prove surjectivity.
\end{proof}

The following formulas will be useful:
\begin{equation}
 S=p^2+up_r>0,\qquad
 z=\frac{\phi p_r-p\phi_r}{S},\qquad
 g=\frac{p\phi+u\phi_r}{S}.
\end{equation}
The Bessel equations give the identities
\begin{align}
 A_r&=A^2+AC+1-C^2,& C_r&=C^2,\label{riccati}\\
 f_r&=w+Af=Tz,&w_r&=(A+C)w-\lambda f,
 &w&=-z-Ag.\label{fw}
\end{align}

\subsection{Two auxiliary functions}
Let
\[
 B(x)=x[J_0(x)^2+J_1(x)^2]-J_0(x)J_1(x).
\]
The recurrence relations give
\[
 B'(x)=\frac{J_0(x)J_1(x)}x.
\]
Define
\begin{equation}\label{Gclosed}
 G(r)=K+\frac{\alpha rB(\alpha r)-KrB(r)}{p(r)^2}.
\end{equation}
The choice of $K$ gives
\[
 B(k)=kc^2,\qquad B(q)=qJ_1(q)^2,\qquad KB(k)=\alpha B(q).
\]
This identity will remove the apparent singularity at $r=k$.

Define
\begin{equation}\label{bclosed}
 b=-KA-GC-fw=-KA+\frac Gr-\frac{\phi\phi_r}{p^2}.
\end{equation}
\begin{lemma}
The function $G$ satisfies
\begin{equation}\label{Videntity}
 V:=w^2-Cfw+\lambda f^2=K A_r+GC^2.
\end{equation}
One also has
\begin{align}
 G_r&=b+2AG+2KC,\label{Gode}\\
 b_r+CG_r&=-2f_rw.\label{bidentity}
\end{align}
\end{lemma}
\begin{proof}
Expand \eqref{Gclosed} using the definition of $B$, multiply by $p^2$,
and use \eqref{riccati}. The result is exactly
$G=r^2(V-KA_r)$, proving \eqref{Videntity}.
By \eqref{fw} and $C_r=C^2$,
\[
 V_r=(2A+C)V-C^2fw.
\]
Also
\[
 A_{rr}=(2A+C)A_r+AC^2-2C^3.
\]
Differentiate \eqref{Videntity}, substitute these two formulas, and
cancel $K(2A+C)A_r$. Division by $C^2>0$ gives
\[
 G_r=(2A-C)G-KA+2KC-fw=b+2AG+2KC.
\]
Differentiating \eqref{bclosed} gives
\[
 b_r+CG_r=-KA_r-GC^2-f_rw-fw_r
         =-V-f_rw-fw_r=-2f_rw,
\]
where the last equality uses
$f_rw-fw_r=w^2-Cfw+\lambda f^2=V$.
\end{proof}

\subsection{Reduction to a two-dimensional quadratic form}
Define the coefficients
\[
 a=K-g^2,\quad d=2G-tG_t,\quad
 L_A=2b-tb_t-2zg,\quad L_C=b+G_t,\quad
 H=-K-z^2+\lambda f^2+b_t.
\]
Define the symmetric matrix
\[
 M=\begin{pmatrix}
 a&d/2&L_A/2\\ d/2&K&L_C/2\\L_A/2&L_C/2&H
 \end{pmatrix}.
\]
\begin{proposition}\label{kernel}
For all $0<r<k$,
\[
 L_C=-dA-2KC,\qquad L_A=-2aA-dC,\qquad
 H=aA^2+dAC+KC^2.
\]

\end{proposition}
\begin{proof}
The first equation follows from \eqref{Gode}:
\[
 dA+2KC+L_C=2GA+2KC+b+(1-tA)G_t
                    =G_r-TG_t=0.
\]
For the second, multiply its left side by $T$ and use
\eqref{bclosed}, $f-g=tz$, and $w=-z-Ag$:
\begin{align*}
 T(L_A+2aA+dC)
 &=T[2b-2zg+2(K-g^2)A+2GC]-t(b_r+CG_r)\\
 &=-2tTzw-t(b_r+CG_r)=0,
\end{align*}
by \eqref{bidentity} and $Tz=f_r$.
Finally \eqref{bidentity} implies $b_t=-2zw-CG_t$. Therefore
\begin{align*}
 H-(aA^2+dAC+KC^2)
 &=\lambda f^2+w^2-K(1+A^2+C^2)-2GAC+CG_r\\
 &=\lambda f^2+w^2-K(1+A^2-C^2)+Cb\\
 &=V-GC^2-K(A^2+AC+1-C^2)=0
\end{align*}
by \eqref{Gode}, \eqref{bclosed}, \eqref{riccati} and
\eqref{Videntity}.
\end{proof}

\begin{proposition}\label{prop:reduction}
If
\begin{equation}\label{remaining-positivity}
 d^2+4Kg^2\le4K^2,
\end{equation}
then $M$ is positive semidefinite and $2K-g^2-d\ge0$.
\end{proposition}
\begin{proof}
The kernel identities give the quadratic-form identity
\[
 (X,Y,P)M(X,Y,P)^T
 =a(X-AP)^2+d(X-AP)(Y-CP)+K(Y-CP)^2.
\]
The block determinant is $K(K-g^2)-d^2/4$; nonnegativity implies
$g^2\le K$. Moreover
\[
 2K-g^2-d\ge2K-g^2-2\sqrt{K(K-g^2)}
            =\left(\sqrt K-\sqrt{K-g^2}\right)^2\ge0.
\]
\end{proof}

\subsection{Analytic endpoint data}\label{sec:endpoints}
Write $\beta=\sqrt K$, $a_0=1-\lambda$ and $\delta=k-r$.
The Bessel equations give
\[
 p=c\delta+\frac{c}{2k}\delta^2+
       \frac{c(3/k^2-1)}6\delta^3+O(\delta^4),\qquad
 u=\frac c2\delta^2+\frac{c}{6k}\delta^3+O(\delta^4),
\]
\[
 \phi=c\beta\delta+\frac{c\beta}{2k}\delta^2+
       \frac{c\beta(2/k^2-\lambda)}6\delta^3+O(\delta^4).
\]
Thus $t=\delta/2+O(\delta^2)$ and, writing $m_k=1-\lambda-k^{-2}$,
\[
 g=\beta-\frac{\beta m_k}{6}\delta^2+O(\delta^3),\qquad
 z=\frac{2\beta m_k}{3}\delta+O(\delta^2).
\]
Here $m_k>0$ follows already from $\lambda<2/5$ and $k>19/5$,
whose elementary rational justifications are supplied below.

The endpoint values of $G$ and its slope are exact. Put
$N(r)=\alpha rB(\alpha r)-KrB(r)$. Differentiating
$(sB(s))'=s(J_0(s)^2+J_1(s)^2)$ gives
\[
 N(k)=N'(k)=0,\quad N''(k)=-2Kc^2,
 \quad N_{rrr}(k)=2Kc^2/k.
\]
Consequently
\[
 G=\frac{Kp^2+N}{p^2}=\frac{2K}{3k}\delta+O(\delta^2),
 \qquad G(k)=0,\quad G_r(k)=-\frac{2K}{3k}.
\]
All these expansions are convergent local analytic expansions after
the indicated removable factors have been cancelled. At the radial center,
$G|_{r=0}=2\lambda-K$, $g|_{r=0}=0$, $z|_{r=0}=(1+c)^{-1}$; these follow
from the ordinary Bessel series.
\section{Convexity of the transform}\label{cv:sec:convexity}
Put $r_0=(1-\lambda)^{-1/2}$. Only the coarse elementary parameter bounds
\begin{equation}
 0<\lambda<\frac25,\qquad 0<c<\frac{41}{100},\qquad k>\frac{19}5
\end{equation}
will be used. They follow from Lemma~\ref{gm:coarse} in
Appendix~\ref{app:constants}. The same appendix verifies
\begin{equation}\label{cv:nineteen}
 \frac7{25}<J_0(19/10)<\frac{283}{1000},\qquad
 \frac{J_1(19/10)}{19/10}>\frac3{10}.
\end{equation}
These use only the degree-six and degree-eight alternating
Bessel bounds, recorded explicitly in the appendix.

\begin{lemma}
One has $z>0$ and $g>0$ throughout $0<r<k$.
\end{lemma}
\begin{proof}
Set $W=p\phi'-\phi p'$. The two Bessel equations give
\[
 (rW)'=r(a_0-r^{-2})p\phi.
\]
Both endpoint limits of $rW$ are zero. Its derivative is negative
before $r_0=a_0^{-1/2}$ and positive after $r_0$. Hence $W<0$ in
the open interval. Thus $f_r=W/p^2<0$ and $z=f_r/t'>0$.

For the other sign put $\psi=-\phi'=\sqrt\lambda
J_1(\sqrt\lambda r)>0$. Now
\[
 [r(p\psi'-\psi p')]'=a_0rp\psi>0.
\]
The expression on the left before differentiation has limit zero
at the center. Consequently $h=\psi/p$ is strictly increasing.
Using $\phi(r)=\int_r^k\psi(s)\dd s$ gives
\[
 p\phi+u\phi'
 =p(r)\int_r^k p(s)[h(s)-h(r)]\dd s>0.
\]
Since $g=(p\phi+u\phi')/S$, the claim follows.
\end{proof}

Put $m=a_0-r^{-2}$ and $Q=a_0u-c$.
\begin{lemma}
For $0<r<k$,
\begin{equation}\label{cv:system}
 g_r=\frac{u}{S}(pmg+Qz),\qquad z_r=-\frac pu g_r.
\end{equation}
\end{lemma}
\begin{proof}
Use $u'=-p$, $p''=-p'/r-(1-r^{-2})p$ and
$\phi''=-\phi'/r-\lambda\phi$. First,
\[
 S'+S/r=p(u/r^2-c).
\]
Differentiating $z=(\phi p'-p\phi')/S$ in $r$ and then dividing by
$t'=-S/p^2$ gives
\[
 f_{tt}=\frac{p^3}{S^3}
 \left[\phi(p'^2-pp'')+\phi''S-\phi'S'\right].
\]
After inserting $p'=u-c-p/r$, the bracket equals
$S[a_0\phi-cz-(p/r^2)g]$. Hence
\[
 f_{tt}=\frac{p^3}{S^2}
       \left[a_0\phi-cz-\frac p{r^2}g\right].
\]
Now use $g_r=-tz_r$, $z_r=f_{tt}t'$ and $\phi=pg+uz$ to obtain
\eqref{cv:system}.
\end{proof}

\begin{lemma}\label{cv:ratio}
There is a unique $r_Q\in(19/10,k)$ for which $Q(r_Q)=0$.
For
\[
 F(r)=\frac{p(r)m(r)}{Q(r)}
\]
one has $F'(r)>0$ both on $(0,r_0)$ and on $(r_Q,k)$.
\end{lemma}
\begin{proof}
Since $Q'=-a_0p<0$, $Q(k)=-c<0$, and
\[
 Q(19/10)>\frac35\frac7{25}-\frac25\frac{41}{100}
 =\frac1{250}>0,
\]
the zero is unique and lies after $19/10$. Also
$r_0\le\sqrt{5/3}<19/10$. Write
$Q_0=Q(0)=a_0-\lambda c>2/5$.

For $0<r<r_0$, put $x=a_0r^2\in(0,1)$. The elementary Bessel
bounds $p\le r/2$ and $p'/p\le1/r$ hold there; the latter follows
from $p'=p/r-J_2(r)$ and the direct alternating-series bound
$J_2(r)\ge r^2/8-r^4/96>0$ for $r<r_0\le\sqrt{5/3}$.
Furthermore
\[
 Q=Q_0-a_0\int_0^r p(s)\dd s\ge Q_0-x/4>0.
\]
As $F<0$, the condition $F'>0$ is equivalent to
\[
 \frac2{r(1-x)}-\frac{p'}p>\frac{a_0p}{Q}.
\]
The left side is at least $(1+x)/[r(1-x)]$, while the right side
is at most $x/[2r(Q_0-x/4)]$. The required strict comparison follows
from
\[
 2(Q_0-x/4)(1+x)-x(1-x)
 =2Q_0+(2Q_0-3/2)x+x^2/2
 >4/5-7x/10+x^2/2\ge1/10.
\]

For $r>r_Q$ put $A=-p'/p$ and $\ell=r-r_Q$. The estimates
\eqref{cv:nineteen} give $p'(19/10)<283/1000-3/10<0$.
Since $(\log p)''<0$ for $r>\sqrt5/2$, we have $A>0$ on
$[r_Q,k)$. The tangent-line bound for the concave function $\log p$
gives, for $r_Q\le s\le r$,
\[
 p(s)\le p(r)e^{A(r)(r-s)},\qquad
 \frac{a_0p(r)}{-Q(r)}
 =\frac{p(r)}{\int_{r_Q}^r p(s)\dd s}
 \ge\frac{A(r)}{e^{A(r)\ell}-1}.
\]
It follows that
\[
 A+\frac{a_0p}{-Q}\ge\frac{A}{1-e^{-A\ell}}>\frac1\ell.
\]
On the other hand,
\[
 \frac{m'}m=\frac2{r(a_0r^2-1)}<\frac1{r-r_Q}.
\]
Indeed, the last strict inequality is equivalent to
$a_0r^3-3r+2r_Q>0$; it holds at $r=r_Q>r_0$ and its derivative is
positive thereafter. Therefore
\[
 \frac\dd{\dd r}\log(-F)
 =-A+\frac{m'}m-\frac{a_0p}{-Q}<0.
\]
As $F<0$ on this interval too, this proves $F'>0$.
\end{proof}

\begin{proposition}
For every $0<r<k$,
\[
 g_r>0,\qquad f_{tt}>0,\qquad 0<g(r)<\sqrt K.
\]
\end{proposition}
\begin{proof}
At any zero of $g_r$ in an interval with $Q\ne0$, the system
\eqref{cv:system} gives $z_r=0$ and
\begin{equation}\label{cv:zeroorientation}
 g_{rr}=\frac{u}{S}\,gQ\left(\frac{pm}{Q}\right)'.
\end{equation}
The endpoint expansions give
\[
 g_r\longrightarrow\frac{a_0-\lambda c}{1+c}>0\quad(r\downarrow0),
\]
and, with $\delta=k-r$,
\[
 g=\beta-\frac{f_2}4\delta^2+O(\delta^3),\qquad
 f_2=\frac{2\beta}{3}(a_0-k^{-2})>0.
\]
Thus $g_r>0$ near both endpoints. On $(0,r_0)$, Lemma~\ref{cv:ratio}
and \eqref{cv:zeroorientation} make every hypothetical zero of $g_r$
a strict crossing from negative to positive. Since $g_r$ starts
positive, no such zero is possible. On $[r_0,r_Q]$ both terms in
$pmg+Qz$ are nonnegative and at least one is positive, so $g_r>0$
directly. On $(r_Q,k)$ every hypothetical zero is instead a strict
crossing from positive to negative. Positivity near $k$ excludes
such a zero. This proves $g_r>0$ globally.

Because $g_r=-t t'f_{tt}$ and $t>0$, $t'<0$, strict convexity
follows. Finally $g(0+)=0$ and $g(k-)=\beta$ give the claimed
bound.
\end{proof}

\section{A quadratic estimate}
The elementary coarse bounds needed below are
\begin{equation}
 \frac14<\lambda<\frac25,\qquad
 \frac25<c<\frac{41}{100},\qquad 2<k<4.
\end{equation}
These are included in Lemma~\ref{gm:coarse}; in particular, the lower
bound on $\lambda$ follows from $q>2$ and $k<4$.

\subsection{Elementary bounds for the phase}
Set
\[
 \rho=\sqrt{u^2+p^2},\qquad \theta=\arctan(p/u),\qquad
 \eta=\theta'=\frac{S}{\rho^2}>0.
\]
The phase increases from zero to $\pi/2$. Write $s_\theta=\sin\theta=p/\rho$
and $c_\theta=\cos\theta=u/\rho$. The system
$u'=-p$, $p'=u-c-p/r$ gives
\begin{equation}\label{ph:phase}
 \eta=1-\frac{c}{\rho} c_\theta-\frac{s_\theta c_\theta}{r},\qquad
 \frac{\rho'}\rho=-\frac{c}{\rho} s_\theta-\frac{s_\theta^2}{r}.
\end{equation}

\begin{lemma}\label{ph:diskbounds}
Throughout $0<r<k$,
\begin{equation}\label{ph:hbounds}
 \eta>\frac{s_\theta c_\theta}{r},\qquad \eta<\frac{\tan\theta}{r},\qquad
 \eta>\frac12-\frac{s_\theta c_\theta}{r}.
\end{equation}
\end{lemma}
\begin{proof}
First $J_2>0$ and $J_3>0$ on $(0,4)$. For $J_3$, the first two
alternating terms give $J_3(r)\ge r^3/48-r^5/768>0$.
For $J_2$, four terms give
\[
 J_2(r)\ge r^2\left(\frac18-\frac{r^2}{96}
                 +\frac{r^4}{3072}-\frac{r^6}{184320}\right)>0.
\]
Indeed the polynomial in parentheses, as a function of $x=r^2$,
is decreasing on $[0,16]$ and has value $7/360$ at $16$.

Let $F=p^2-uJ_2$. The recurrences $p'=p/r-J_2$ and
$J_2'=p-2J_2/r$ give the exact first-order equation
\[
 (r^2F)'=rp(2p-cr).
\]
Since $(p/r)'=-J_2/r<0$, the forcing has exactly one change of sign,
from positive to negative. Both endpoint limits of $r^2F$ are zero.
Consequently $F>0$ in the open interval. It follows that
$S=up/r+F>up/r$, proving the first bound in \eqref{ph:hbounds}.

The recurrence $4J_2/r=p+J_3$ gives another exact identity:
\[
 p \rho^2-ruS=p\left(p-\frac{ru}2\right)^2
                  +\frac{r^2u^2}{4}J_3>0.
\]
Division by $ru\rho^2$ proves the second bound.

Finally $(J_0^2+J_1^2)'=-2p^2/r<0$ and
$J_0(k)^2+J_1(k)^2=c^2$. Hence
\[
 \rho^2-2cu=J_0(r)^2+p(r)^2-c^2>0.
\]
Thus $cu/\rho^2<1/2$, which combined with \eqref{ph:phase} gives the
third bound.
\end{proof}

\begin{lemma}\label{ph:radius}
If $\theta\ge\pi/4$, then $r>2$ and
\[
 r>2\sqrt2\sin\theta.
\]
\end{lemma}
\begin{proof}
The elementary series bounds $J_0(2)>11/50$,
$J_1(2)<3/5$, together with $c>2/5$, give $u(2)>p(2)$.
Thus $\theta(2)<\pi/4$ and $\theta\ge\pi/4$ implies $r>2$.

By the second identity in the preceding proof,
\begin{equation}\label{ph:logsine}
 \left(\log\frac{\sin\theta}{r}\right)'
 =\cot\theta\left(\eta-\frac{\tan\theta}{r}\right)
 \le-\frac{ru^2J_3}{4p \rho^2}<0.
\end{equation}
On $0<r\le2$ we have $u^2/\rho^2>1/2$, $p\le r/2$ and
\[
 J_3(r)\ge\frac{r^3}{48}-\frac{r^5}{768}
          \ge\frac{r^3}{64}.
\]
Therefore the right side of \eqref{ph:logsine} is at most
$-r^3/256$. Integrating from zero to $2$, using
$\lim_{r\downarrow0}\sin\theta/r=[2(1+c)]^{-1}\le5/14$,
gives
\[
 \frac{\sin\theta(2)}2
 <\frac5{14}e^{-1/64}<\frac5{14}\frac{64}{65}
 =\frac{32}{91}<\frac1{2\sqrt2}.
\]
The last inequality is the integer comparison $8192<8281$.
Monotonicity in \eqref{ph:logsine} extends this bound to every $r\ge2$.
\end{proof}

\subsection{Monotonicity of an auxiliary potential}
\begin{proposition}\label{ph:potential}
The function
\[
 \mathcal V(r)=-\frac{\rho''+\rho'/r}{\rho}
\]
is strictly decreasing on $(0,k)$.
\end{proposition}
\begin{proof}
Direct differentiation of \eqref{ph:phase} gives
\[
 \eta'=\frac{s_\theta c_\theta}{r^2}+(2\eta-1)\frac{c}{\rho} s_\theta
                     +\frac{\eta(2s_\theta^2-1)}r,
\]
\[
 \mathcal V=1-\frac{c}{\rho} c_\theta-\frac{s_\theta^2}{r^2}-\eta^2,
\]
and consequently
\begin{equation}\label{ph:Vprime}
 \mathcal V'=-\frac{cp}{\rho^2}(1-2\eta)^2
 +\frac2r\left[\left(\eta c_\theta-\frac{s_\theta}{r}\right)^2-\eta^2s_\theta^2\right].
\end{equation}
We show that the square bracket is strictly negative.

For $0<\theta<\pi/4$ the bracket, as a quadratic in $\eta$, has
positive roots
\[
 \eta_- =\frac{s_\theta}{r(s_\theta+c_\theta)},\qquad \eta_+=\frac{s_\theta}{r(c_\theta-s_\theta)}.
\]
Lemma~\ref{ph:diskbounds} gives
\[
 \eta>\frac{s_\theta c_\theta}{r}\ge \eta_-,\qquad
 \eta<\frac{s_\theta}{r c_\theta}<\eta_+.
\]
Here $c_\theta(s_\theta+c_\theta)\ge1$ is equivalent to $c_\theta\ge s_\theta$.
At $\theta=\pi/4$ the bracket is linear in $\eta$ and the same lower
bound still makes it strictly negative.

For $\pi/4<\theta<\pi/2$ the quadratic has negative leading
coefficient and is negative whenever $\eta>\eta_-$. By
Lemma~\ref{ph:radius}, $r>2\sqrt2s_\theta$. On this angular range
\[
 c_\theta+\frac1{s_\theta+c_\theta}\le\sqrt2.
\]
Indeed its value at $\pi/4$ is $\sqrt2$, and its derivative is
$-s_\theta+(s_\theta-c_\theta)/(s_\theta+c_\theta)^2<0$. It follows that
\[
 r>2s_\theta\left(c_\theta+\frac1{s_\theta+c_\theta}\right),\qquad
 \frac12-\frac{s_\theta c_\theta}{r}>\frac{s_\theta}{r(s_\theta+c_\theta)}=\eta_-.
\]
The third bound in \eqref{ph:hbounds} again gives $\eta>\eta_-$.
Thus the bracket in \eqref{ph:Vprime} is negative everywhere, proving
$\mathcal V'<0$.
\end{proof}

\subsection{A Wronskian comparison}
\begin{proposition}
For every $t>0$,
\[
 g^2+z^2<K.
\]
\end{proposition}
\begin{proof}
The endpoint expansions give
\[
 \mathcal V(0+)=\frac{c+1/2}{(1+c)^2}>\frac25>\lambda,
 \qquad
 \mathcal V(k-)=\frac14-\frac1{k^2}<\frac14<\lambda.
\]
For the first strict comparison, a sufficient rational bound is
$\mathcal V(0+)>(9/10)/(141/100)^2>2/5$.
By Proposition~\ref{ph:potential}, $\mathcal V-\lambda$ has exactly one sign
change, from positive to negative.

Let $W=\rho\phi'-\phi \rho'$. Since $\phi''+\phi'/r+\lambda\phi=0$,
\[
 (rW)'=r \rho\phi(\mathcal V-\lambda).
\]
Both endpoint limits of $rW$ are zero. The forcing is first positive
and then negative, so $W>0$ throughout $(0,k)$. Using
$\rho'=-p(c+p/r)/\rho$ and the formulas for $z,g$ gives the exact relation
\[
 z-tg=-\frac{\rho W}{pS}<0.
\]
Finally $g_t=-tf_{tt}$ and $z_t=f_{tt}$ imply
\[
 \frac\dd{\dd t}(g^2+z^2)=2f_{tt}(z-tg)<0.
\]
As $t\downarrow0$, $(g,z)\to(\beta,0)$, giving the energy bound.
\end{proof}

\section{A lower bound for the first derivative}
In this section put $v=1-J_0(r)$ and $\sigma=1+c=u+v$.
We first need an estimate for the disk profile.

\begin{lemma}\label{dc:ratio}
On $0<r<k$,
\begin{equation}\label{dc:boundB}
 \frac{2vS}{p^2}<\sigma.
\end{equation}
\end{lemma}

\begin{proof}
The standard Bessel product and derivative recurrence give
\[
 J_1(r)=\frac r2\prod_{n\ge1}
                  \left(1-\frac{r^2}{j_{1,n}^2}\right),\qquad
 \frac{J_2(r)}{rJ_1(r)}
       =\sum_{n\ge1}\frac2{j_{1,n}^2-r^2}.
\]
The series and its derivative converge uniformly on each compact
subinterval of $(0,k)$.  In particular $J_2/(rJ_1)$ is strictly
increasing there.  See \cite[Eqs.~10.21.15 and 10.6.2]{DLMF}.

Since $X=J_0(r)$ is strictly decreasing from $1$ to $-c$, regard
$E(X)=p(r)^2$ as a function of $X\in[-c,1]$.  The function $E$ is
continuous on the closed interval, vanishes at its endpoints and
is smooth in its interior.  Using $X'=-p$ and
$p'=X-p/r$ gives
\[
 E_X=-2X+\frac{2p}r,\qquad
 E_{XX}=-2+\frac{2J_2(r)}{rp},\qquad E_{XXX}<0.
\]
The last sign follows because $J_2/(rp)$ increases with $r$ while
$X$ decreases with $r$.

For fixed endpoints $x_-=-c$, $x_+=1$, the second divided difference is
\[
 E[x_-,X,x_+]=-\frac{E(X)}{(x_+-X)(X-x_-)}.
\]
Its derivative with respect to the middle node is the third divided
difference $E[x_-,X,X,x_+]$.  The Hermite divided-difference mean-value
theorem yields, for some $\xi\in(x_-,x_+)$,
\[
 \frac{d}{dX}E[x_-,X,x_+]=E[x_-,X,X,x_+]
                   =\frac{E_{XXX}(\xi)}{3!}<0.
\]
Only continuity of $E$ at $x_-,x_+$ is needed here: the usual proof
subtracts its cubic Hermite interpolant and applies Rolle's theorem
three times to interior points.  No endpoint second derivative is
being assumed.

It follows that $E(X)/[(1-X)(X+c)]$ is strictly increasing with
$X$, and hence strictly decreasing with $r$.  Differentiating that
ratio gives precisely
\[
 \left(\frac{p^2}{uv}\right)'
 =-\frac{p}{u^2v^2}\{(u-v)p^2-2uvp'\}<0.
\]
Finally $\sigma p^2-2vS=(u-v)p^2-2uvp'$, proving
\eqref{dc:boundB}.
\end{proof}

\begin{corollary}
Put $h_0=(p+vA)/\sigma$. The function
\[
 \mathcal R=\lambda-\frac{h_0}{r}+h_0^2-h_0'
\]
is strictly decreasing on $(0,k)$.
\end{corollary}
\begin{proof}
Logarithmic differentiation of the same Bessel product gives
$A'>0$.  Also $p/r$ is strictly decreasing, so
$rp<2\int_0^rp(s)\,ds=2v$.  Direct differentiation, using
$h_0'=vA'/\sigma$, gives
\[
 \mathcal R'
 =\frac{pA'}{\sigma}\left(\frac{2vS/p^2}{\sigma}-1\right)
                      +\frac{rp-2v}{\sigma r^3}<0
\]
by Lemma~\ref{dc:ratio}.
\end{proof}

We will also use $J_0(qs)>1-s$ for $0<s<1$. Here is an elementary proof.
$J_1$ is strictly unimodal on $(0,q)$: its derivative is positive
for $r\le\sqrt5/2$ by the change-of-variable lemma, and at any later zero its
second derivative is $-(1-r^{-2})J_1$, so a downward crossing
can occur only after $1$, and all such later crossings are downward.
The elementary series gives $J_1'(1)>5/16>0$, and
$J_1'(q)=-J_1(q)/q<0$, so there is exactly one such crossing.
Moreover $qJ_1(q)>1$.
Indeed $2<q<5/2$ and the four-term alternating lower bound gives,
with $s=q^2\in[4,25/4]$,
\[
 qJ_1(q)\ge Q(s),\quad
 Q(s)=\frac s2-\frac{s^2}{16}+\frac{s^3}{384}
                                      -\frac{s^4}{18432}.
\]
Here $Q''(s)=-[(s-12)^2+48]/1536<0$ and
$Q(4)=83/72>1$, $Q(25/4)=5834975/4718592>1$.
Therefore $J_1(r)-1/q$ has exactly one sign change on $(0,q)$,
from negative to positive.  The function $J_0(r)-1+r/q$ first
increases and then decreases, with zero values at both endpoints,
which proves the assertion.

\begin{proposition}
For $0<r<k$,
\[
 z(r)>\frac{k-r}{k(1+c)}.
\]
\end{proposition}
\begin{proof}
Use $h_0$ from the preceding corollary.
Let $\psi$ be the positive solution of
\[
 \psi'=-h_0\psi,\qquad \psi(0)=1.
\]
At the center $h_0=r/(4\sigma)+O(r^3)$, so $\psi$ extends as a smooth
even positive function.  At the other endpoint, with $\delta=k-r$,
the elementary Bessel expansions give
\[
 p=c\delta+\frac{c}{2k}\delta^2+O(\delta^3),\quad
 v=\sigma-\frac c2\delta^2+O(\delta^3),\quad
 h_0=\frac1\delta+\frac1{2k}+O(\delta).
\]
Consequently $\psi=c_\psi\delta(1+\delta/(2k)+O(\delta^2))$ for
some $c_\psi>0$.  In particular $\psi$ and $\phi$ have simple zeros
at $k$.

Define
\[
 V_\psi=-\frac{\psi''+\psi'/r}{\psi}
       =h_0'+h_0/r-h_0^2=\lambda-\mathcal R.
\]
By the preceding corollary $V_\psi$ is strictly increasing.
For $W=\psi\phi'-\phi\psi'$ one has the exact identity
\[
 (rW)'=r\psi\phi(V_\psi-\lambda).
\]
The central regularity and the simple boundary zeros give
$rW\to0$ at both endpoints.  The strictly increasing function
$V_\psi-\lambda$ must therefore change sign exactly once,
from negative to positive: otherwise its integral against the
positive weight $r\psi\phi$ could not be zero.  Thus $rW$
first decreases from zero and then increases to zero, and
$W<0$ throughout the open interval.  It follows that
\[
 -\frac{\phi'}\phi-h_0=-\frac W{\psi\phi}>0.
\]
Finally the exact transformed derivative is
$z=(\phi p'-p\phi')/S$.  Direct subtraction, using
$\sigma-u=v$, yields
\[
 z-\frac{\phi}{\sigma}
 =\frac{\phi p}{S}
       \left(-\frac{\phi'}\phi-\frac{p-vp'/p}{\sigma}\right)
 =\frac{\phi p}{S}
       \left(-\frac{\phi'}\phi-h_0\right)>0.
\]
The inequality for $J_0$ proved above, rescaled to $(0,k)$, gives
$\phi(r)>(k-r)/k$ and completes the proof.
\end{proof}

\section{Concavity of the auxiliary function}\label{sec:G}
Write $q_n=j_{0,n}$ and $\kappa_n=j_{1,n}$.
\begin{lemma}\label{ev:Benvelope}
For every $n\ge1$ and every $x\ge q_n$,
\[
 B(\kappa_n)\le B(x)\le B(q_n).
\]
Equality on the left and right holds only at $x=\kappa_n$ and $x=q_n$,
respectively.
\end{lemma}
\begin{proof}
One has $B'(x)=J_0(x)J_1(x)/x$. By interlacing, the positive zeros
$q_n$ of $J_0$ give local maxima of $B$ and the positive zeros $\kappa_n$
of $J_1$ give local minima. The identities
\[
 2B(x)-\frac{J_1(x)^2}{x}
       =3\int_0^x\frac{J_1(s)^2}{s^2}ds
\]
and
\[
 B(q_{n+1})-B(q_n)
   =-\frac12\int_{q_n}^{q_{n+1}}\frac{J_0(s)^2}{s^2}ds<0
\]
show, respectively, that the successive minima increase strictly and
the successive maxima decrease strictly. Between successive critical
points $B$ is strictly monotone. The asserted bounds follow for each $n$. Also $B(q)>B(k)>0$ and
$K=\alpha B(q)/B(k)$.
\end{proof}

Put $L(r)=KB(r)-\alpha B(\alpha r)$, so that
$G(r)=K-rL(r)/p(r)^2$. The endpoint identities give $G(k)=0$ and
\begin{equation}\label{gm:endslope}
 G_r(k)=-\frac{2K}{3k}.
\end{equation}
The singularities at $0,k,-k$ are removable. The function $G$ is an
even meromorphic function whose only other possible poles are double
poles at $\pm\kappa_n$, $n\ge2$. There are no unlisted complex zeros:
if $J_1(\zeta)=0$ with $\zeta\ne0$, multiply the angular-one radial
Dirichlet equation for $v(s)=J_1(\zeta s)$ on $(0,1)$ by $\bar v$ to get
\[
 \int_0^1\left(s|v'|^2+\frac{|v|^2}s\right)\,ds
 =\zeta^2\int_0^1s|v|^2\,ds.
\]
The two integrals are strictly positive, so $\zeta^2>0$ and
$\zeta$ is real. The nonzero
zeros are simple by uniqueness for the Bessel ODE.

The fixed bounds needed in the following lemma, including
$K>.654$, $B(\kappa_3)>.634$, $B(q_2)<.640$ and
$J_1(q_2)^2<.116$, are proved in Lemma~\ref{gm:coarse}.

\begin{lemma}\label{gm:residuepositive}
For $n\ge2$, $L(\kappa_n)>0$. Moreover
\[
 3L(\kappa_n)-\phi(\kappa_n)\phi_r(\kappa_n)>0.
\]
\end{lemma}
\begin{proof}
Lemma~\ref{ev:Benvelope} gives $B(x)\le B(q)$ for $x\ge q$ and
$B(\kappa_n)>B(k)$ for $n\ge2$. Since $KB(k)=\alpha B(q)$,
this proves $L(\kappa_n)>0$.

For $n=2$, the inequalities $k<\alpha\kappa_2<q_2$ imply
$\phi(\kappa_2)<0$ and $\phi_r(\kappa_2)>0$, so the second
claim is immediate. For $n\ge3$ we have
$B(\kappa_n)\ge B(\kappa_3)>.634$ and
$B(\alpha\kappa_n)\le B(q_2)<.640$. Also
$E(x)=J_0(x)^2+J_1(x)^2$ satisfies $E'=-2J_1(x)^2/x\le0$.
Thus for $x=\alpha\kappa_n\ge q_2$,
\[
 |\phi(\kappa_n)\phi_r(\kappa_n)|
 =\alpha|J_0(x)J_1(x)|\le\frac\alpha2 E(x)
 <\frac{.628\cdot.116}{2}.
\]
The coarse bounds now give the strictly positive rational estimate
\[
 3L(\kappa_n)-\phi(\kappa_n)\phi_r(\kappa_n)
 >3(.654\cdot.634-.628\cdot.640)-\frac{.628\cdot.116}2
 =\frac{431}{250000}>0.
\]
\end{proof}

\subsection{Uniform complex asymptotics and coefficient residues}
For the coefficient calculation we use the classical uniform asymptotic
formula, for $\nu=0,1$,
\begin{equation}\label{asym:uniformJ}
 J_\nu(\zeta)=\left(\frac2{\pi\zeta}\right)^{1/2}
 \left[\cos\left(\zeta-\frac{\nu\pi}2-\frac\pi4\right)
       +O\left(\frac{e^{|\operatorname{Im}\zeta|}}{|\zeta|}\right)\right],
 \qquad |\arg\zeta|\le\frac\pi2.
\end{equation}
The error is uniform on the closed sector, with the principal square
root. This follows from the Hankel-function remainder estimates in
\cite[\S10.17(iv), Eqs.~10.17.13--10.17.15]{DLMF} and
$2J_\nu=H_\nu^{(1)}+H_\nu^{(2)}$; the variation bound is at most
$\pi/(2|\zeta|)$ on the relevant half-plane, and the conjugate bound
handles the other Hankel function. Integer-order parity extends the
absolute-value estimates to the left half-plane.

\begin{lemma}\label{asym:circles}
For $R_N=(N+3/4)\pi$ and all sufficiently large integers $N$,
$G$ has no pole on $|\zeta|=R_N$ and
\[
 \sup_{|\zeta|=R_N}|G(\zeta)|\le C R_N^2,
\]
where $C$ is independent of $N$.
\end{lemma}
\begin{proof}
Write $R=R_N$ and $\zeta=x+iy$. By parity assume $x\ge0$.
If $|y|\ge1$, then
\[
 |\cos(\zeta-3\pi/4)|\ge\sinh|y|\ge c_0 e^{|y|}
\]
for a fixed $c_0>0$. If $|y|<1$, then
$0\le R-x=y^2/(R+x)\le1/R$ and $R-3\pi/4=N\pi$, so
$|\cos(\zeta-3\pi/4)|\ge\cos(1/R)$.
Thus \eqref{asym:uniformJ}, with its error absorbed for large $R$, gives
\[
 |J_1(\zeta)|\ge c_1R^{-1/2}e^{|y|}.
\]
The same asymptotic formula gives
\[
 |J_0(\zeta)|+|J_1(\zeta)|\le C_1R^{-1/2}e^{|y|},\qquad
 |J_0(\alpha\zeta)|+|J_1(\alpha\zeta)|
       \le C_2R^{-1/2}e^{\alpha|y|}.
\]
Here $\alpha$ is fixed and positive. Consequently
$|B(\zeta)|+|B(\alpha\zeta)|\le C_3e^{2|y|}$.
The representation $G=K-\zeta L/J_1^2$ now proves the claim.
\end{proof}

For a function analytic at zero, $[r^j]G$ denotes the coefficient of $r^j$ in its Taylor series.
\begin{lemma}
For every integer $m\ge2$,
\begin{equation}\label{gm:resformula}
 -[r^{2m}]G(r)=\sum_{n=2}^{\infty}
 \frac{2[(2m-1)L(\kappa_n)-\phi(\kappa_n)\phi_r(\kappa_n)]}
 {\kappa_n^{2m+1}J_0(\kappa_n)^2}.
\end{equation}
The series is absolutely convergent.
\end{lemma}
\begin{proof}
At $\kappa=\kappa_n$ let $x_\kappa=J_0(\kappa)\ne0$.
The Bessel equation gives
\[
 p(r)=x_\kappa(r-\kappa)
       -\frac{x_\kappa}{2\kappa}(r-\kappa)^2
       +O((r-\kappa)^3).
\]
Therefore the principal part of $K-G=rL/p^2$ at $\kappa$ is
\[
 \frac{A_\kappa}{(r-\kappa)^2}
    +\frac{D_\kappa}{r-\kappa},\qquad
 A_\kappa=\frac{\kappa L(\kappa)}{x_\kappa^2},\quad
 D_\kappa=\frac{2L(\kappa)+\kappa L'(\kappa)}{x_\kappa^2}.
\]
Since $B'(\kappa)=0$,
$\kappa L'(\kappa)=\phi(\kappa)\phi_r(\kappa)$.
Evenness determines the principal part at $-\kappa$.

To justify the coefficient identity, apply the residue theorem to
$G(\zeta)/\zeta^{2m+1}$ on $|\zeta|=R_N$, using Lemma~\ref{asym:circles}.
For each fixed integer $m\ge2$,
\begin{equation}\label{asym:outervanishes}
 \left|\frac1{2\pi i}\int_{|\zeta|=R_N}\frac{G(\zeta)}{\zeta^{2m+1}}\dd\zeta\right|
       \le C_{\alpha,K}R_N^{2-2m}\longrightarrow0.
\end{equation}
The residue at zero is $[\zeta^{2m}]G(\zeta)$. With the principal parts already
computed, the sum of the residues at $\kappa$ and $-\kappa$ is
\[
 \frac{2(2m+1)A_\kappa}{\kappa^{2m+2}}
       -\frac{2D_\kappa}{\kappa^{2m+1}}
 =\frac{2[(2m-1)L(\kappa)-\phi(\kappa)\phi_r(\kappa)]}
        {\kappa^{2m+1}J_0(\kappa)^2}.
\]
It remains to verify absolute convergence before taking the limit of
the finite residue sums. The zero asymptotic
\cite[Eq.~10.21.19]{DLMF}
gives
\[
 \kappa_n=(n+1/4)\pi+O(n^{-1}).
\]
Apply \eqref{asym:uniformJ} on the positive axis at $\kappa_n$ to obtain
\[
 J_0(\kappa_n)=(-1)^n\left(\frac{2}{\pi\kappa_n}\right)^{1/2}
                         (1+O(\kappa_n^{-1})).
\]
The same upper bounds give $B(x)=O(1)$ as $x\to+\infty$, so
$L(\kappa_n)=O(1)$ and
\[
 \phi(\kappa_n)\phi_r(\kappa_n)
   =-\alpha J_0(\alpha\kappa_n)J_1(\alpha\kappa_n)
   =O(\kappa_n^{-1}).
\]
Consequently, for each fixed $m\ge2$, the absolute value of the
displayed paired residue is $O_m(\kappa_n^{-2m})$, and the resulting
series converges absolutely. Letting $N\to\infty$ in the residue
theorem and using \eqref{asym:outervanishes} gives
\[
 -[\zeta^{2m}]G(\zeta)=\sum_{n=2}^{\infty}
 \frac{2[(2m-1)L(\kappa_n)-\phi(\kappa_n)\phi_r(\kappa_n)]}
        {\kappa_n^{2m+1}J_0(\kappa_n)^2}.
\]

\end{proof}

\subsection{Concavity and the final estimates}
\begin{proposition}
There are numbers $a_m>0$ such that
\[
 G(r)=2\lambda-K-\sum_{m=1}^\infty a_m r^{2m},
 \qquad |r|<\kappa_2.
\]
For $0<r<k$,
\[
 G(r)>0,\qquad G_r(r)<0,\qquad G_{rr}(r)<0,
 \qquad G_t(t)>0.
\]
\end{proposition}
\begin{proof}
Lemma~\ref{gm:residuepositive} makes every summand of
\eqref{gm:resformula} positive for $m\ge2$. For $m=1$, direct expansion
at zero gives
\[
 [r^2]G=\frac{2\lambda-\lambda^2-K}{4}<0,
\]
because $0<\lambda<.4$ and $K>.654>.64$.
There are no singularities in $|r|<\kappa_2$, so its Taylor series
converges there. Termwise differentiation proves both strict
derivative signs. Since $G(k)=0$, positivity follows. Finally
$t_r<0$ implies $G_t=G_r/t_r>0$.
\end{proof}

\begin{corollary}
For $\delta=k-r$, $0\le r\le k$,
\[
 0\le G(r)\le\frac{2K}{3k}\delta,\qquad
 0\le-G_r(r)\le\frac{2K}{3k}.
\]
\end{corollary}
\begin{proof}
Strict concavity and $G_r(0)=0$ show that $G_r$ decreases to the
value in \eqref{gm:endslope}. A concave function lies below its tangent
line at $k$, where $G(k)=0$.
\end{proof}

\begin{lemma}\label{gm:tailratio}
For $0<r<k$,
\[
 0<\frac{up}{S}<\frac{k^2-r^2}{r}.
\]
\end{lemma}
\begin{proof}
Write $b=k^2$, $x=r^2$, $p=rh(x)$ and use the canonical product
\[
 h(x)=\frac12(1-x/b)P(x),\qquad
 P(x)=\prod_{n\ge2}(1-x/\kappa_n^2).
\]
Its logarithmic derivative is
\[
 \frac{h'}h=-\frac1{b-x}-\Sigma(x),\qquad
 \Sigma(x)=\sum_{n\ge2}\frac1{\kappa_n^2-x}\le\Sigma(b)=\frac1b.
\]
The final equality follows without a sum evaluation: the endpoint ODE
$p''(k)/p'(k)=-1/k$, applied to the displayed factored product,
gives $P'(b)/P(b)=-1/b$.
Since $P$ decreases, $h(y)\le h(x)(b-y)/(b-x)$ for $x\le y\le b$.
Consequently
\[
 u=\frac12\int_x^bh(y)dy\le\frac{b-x}{4}h(x),\qquad
 S=xh^2+u(h+2xh').
\]
The claimed upper bound is equivalent to positivity of
\[
 J=(b-x)xh^2+uD,\qquad D=(b-2x)h+2x(b-x)h'.
\]
If $D\ge0$ this is immediate. If $D<0$, the upper bound on $u$
gives
\[
 J\ge\frac{(b-x)h^2}{4}
       [b-2x(b-x)\Sigma(x)]
 \ge\frac{(b-x)h^2}{4}
       \left[b-\frac{2x(b-x)}b\right]>0,
\]
since $x(b-x)\le b^2/4$.
\end{proof}

\begin{corollary}
On $0<r<k$,
\[
 0<d(r)<2G(r)\le\frac{4K}{3k}(k-r).
\]
\end{corollary}
\begin{proof}
Since $t_r=-S/p^2$, $d=2G+C_{\rm rad} G_r$ with $C_{\rm rad}=up/S$.
The positive-coefficient expansion and $G(k)=0$ give
\[
 d=2\sum_{m\ge1}a_m
       [(k^{2m}-r^{2m})-mC_{\rm rad} r^{2m-1}].
\]
Lemma~\ref{gm:tailratio} shows that each bracket is strictly positive,
because
\[
 k^{2m}-r^{2m}
 =(k^2-r^2)\sum_{j=0}^{m-1}k^{2(m-1-j)}r^{2j}
 \ge m(k^2-r^2)r^{2m-2}>mC_{\rm rad} r^{2m-1}.
\]
This proves $d>0$. The strict upper bound $d<2G$ follows from
$tG_t>0$, and the final linear estimate was proved above.
\end{proof}

\subsection{Completion of the matrix estimate}
The preceding sections give, on $0<r<k$,
\[
 g^2+z^2<K,\qquad
 0<d<2G\le\frac{4K}{3k}(k-r),\qquad
 z>\frac{k-r}{k(1+c)}.
\]
Since $G|_{r=0}=2\lambda-K>0$, one has $K<2\lambda<4/5<1$, and $c<1/2$.
It follows that $2\sqrt K(1+c)<3$. Therefore
\[
 0<d<\frac{4K}{3k}(k-r)
 <2\sqrt K\frac{k-r}{k(1+c)}<2\sqrt K\,z.
\]
Together with the quadratic bound for $g$ and $z$ this proves
\begin{equation}\label{analytic-coupling}
 \boxed{\quad d^2<4Kz^2<4K(K-g^2),\qquad0<r<k.\quad}
\end{equation}
By \eqref{analytic-coupling}, the matrix $M$ and the separate mixed-Hessian
coefficient are positive semidefinite by Proposition~\ref{kernel}
and Proposition~\ref{prop:reduction}. In particular,
\eqref{remaining-positivity} holds throughout the parameter interval.
\section{The functional inequality}\label{sec:functional}
We first record the endpoint control needed to apply the profiles to
arbitrary smooth functions, including those that change sign. We write
$\|\cdot\|_2$ for the $L^2$ norm on the domain under consideration
and use the Frobenius norm for matrices.

\begin{lemma}\label{lem:flux-endpoints}
With
$\ell=(1+c)^{-1}$ and $\beta=\sqrt K$, their endpoint behavior is
\begin{align*}
 f(t)&=\beta+O(t^2),& b(t)&=O(t),&
 G(t)&=\frac{4K}{3k}t+O(t^2) &&(t\downarrow0),\\
 f(t)&=\ell t+O(t^{-1}),&
 b(t)&=2\lambda\ell t+O(t^{-1}),&
 G(t)&=2\lambda-K+O(t^{-2}) &&(t\to\infty).
\end{align*}
The expansions are convergent analytic or Laurent expansions at the
respective endpoints. The functions
\[
 f_t,\quad f-tf_t,\quad b_t,\quad b-tb_t,
 \quad G,\quad G_t,\quad tG_t
\]
are bounded on $(0,\infty)$.
\end{lemma}
\begin{proof}
At $r=k$, the expansions in Section~\ref{sec:endpoints} show that
$t=\delta/2+O(\delta^2)$ has an analytic inverse and
$f=\beta+O(\delta^2)$. The formula for $G$ gives the stated expansion.
Moreover
\[
 A=\delta^{-1}+\frac1{2k}+O(\delta),\qquad
 fw=-K\delta^{-1}-\frac K{2k}+O(\delta).
\]
In $b=-KA+G/r-fw$, both the pole and constant terms cancel. Hence
$b=O(\delta)$, and all three functions are analytic in $t$ at zero.

At $r=0$, the function $y=1/t=p/u$ is odd and analytic, with
$y_r(0)=1/[2(1+c)]$. Its inverse $r=r(y)$ is odd and analytic.
The functions $f=\phi/p$ and $b$ are odd meromorphic functions of $r$
with at most a simple pole, while $G$ is even and analytic there.
The leading terms are
\[
 t=\frac{2(1+c)}r+O(r),\qquad
 f=\frac2r+O(r),\qquad b=\frac{4\lambda}r+O(r),\qquad
 G=2\lambda-K+O(r^2).
\]
Substitution of $r=r(1/t)$ proves the expansions at infinity and the
claim about differentiated remainders. The boundedness assertions
follow at the endpoints from these expansions and on compact
subintervals from smoothness.
\end{proof}

Extend $f$ and $G$ evenly and $b$ oddly to the real line. The even
extension of $G$ is Lipschitz, not differentiable at zero; only its
first weak derivative will be used. Put
\[
 \mathcal F(a,P)=
 \begin{cases}P f(a/P),&P>0,\\ \ell|a|,&P=0,\end{cases}
 \qquad \ell=(1+c)^{-1}.
\]
The following is the functional form of the comparison.
\begin{proposition}\label{prop:functional}
For every real $U\in C_c^\infty(\mathbb R^2)$, the function
$F=\mathcal F(U,|\nabla U|)$ is compactly supported and Lipschitz, and
\begin{equation}\label{functional}
 \int_{\mathbb R^2}|\nabla F|^2-\lambda\int_{\mathbb R^2}F^2
 \le K\left(\int_{\mathbb R^2}|\Delta U|^2
                     -\int_{\mathbb R^2}|\nabla U|^2\right).
\end{equation}
\end{proposition}
\begin{proof}
On $\{P:=|\nabla U|>0\}$ choose the orthonormal frame
$e_1=\nabla U/P$, $e_2=(-\partial_2U,\partial_1U)/P$, and write
\[
 X=e_1^TD^2Ue_1,\qquad Y=e_2^TD^2Ue_2,
 \qquad Z=e_1^TD^2Ue_2,\qquad s=U/P.
\]
In the following identities all profiles are evaluated at $s$.
The chain rule gives
\[
 |\nabla F|^2=(zP+gX)^2+g^2Z^2.
\]
For a $2\times2$ matrix, $\cof$ denotes its cofactor matrix. Direct
calculation, or the identity $\operatorname{div}(\cof D^2U)=0$, gives
\begin{align*}
 \operatorname{div}\bigl(Pb(s)\nabla U\bigr)
   &=P\bigl((2b-sb_s)X+bY\bigr)+b_sP^2,\\
 \operatorname{div}\bigl(G(s)\cof(D^2U)\nabla U\bigr)
   &=(2G-sG_s)(XY-Z^2)+PG_sY.
\end{align*}
Adding these two expressions to
\[
 K(|D^2U|^2-P^2)-(|\nabla F|^2-\lambda F^2)
\]
produces exactly
\[
 (X,Y,P)M(s)(X,Y,P)^T+(2K-g^2-d)Z^2.
\]
For $s>0$ this is nonnegative by the matrix estimate proved in
Section~\ref{sec:G}. For $s<0$, the parity extensions give
\[
 M(-s)=D M(s)D,\qquad D=\operatorname{diag}(-1,-1,1),
\]
and the same conclusion follows.

It remains to justify the calculation across the nodal and critical
sets. If a scalar profile $a$ has bounded $a_s$ and $a-sa_s$, then
$P a(A/P)$ has bounded derivatives in $(A,P)$ for $P>0$. Integrating
these bounds separately in $A$ and $P$ gives a Lipschitz extension to
$\mathbb R\times[0,\infty)$. Lemma~\ref{lem:flux-endpoints} applies
this observation to $f$ and $b$. Their extensions at $P=0$ are
$\ell|A|$ and $2\lambda\ell A$, respectively. Therefore $F$ and
\[
 V_1=P b(U/P)\nabla U
\]
with this extension are compactly supported Lipschitz functions.

For the second vector field, define
\[
 \Psi(A,\xi)=G(A/|\xi|)\xi\quad(\xi\ne0),\qquad
 \Psi(A,0)=0.
\]
Its derivative in $A$ has norm at most $\|G_s\|_\infty$ and its
matrix derivative in $\xi$ is
\[
 G I-sG_s\frac{\xi\otimes\xi}{|\xi|^2},\qquad s=A/|\xi|.
\]
The bounds on $G,G_s,sG_s$ prove that $\Psi$ is globally Lipschitz;
continuity at $\xi=0$ follows from boundedness of $G$. Thus
\[
 V_2=\cof(D^2U)\Psi(U,\nabla U)
\]
is also compactly supported and Lipschitz. Both divergences are
integrable weak derivatives and have integral zero. In particular,
the cusp in $G$ introduces no measure supported on $\{U=0\}$.

The portion of $\{U=0\}$ where $P>0$ is locally a smooth curve and
has area zero. On $\{P=0\}$ one has $D^2U=0$ almost everywhere:
apply the implicit function theorem to each first derivative of
$U$ to discard the zero sets on which its gradient is nonzero.
At each remaining critical point, the derivative of
$x\mapsto(U(x),\nabla U(x))$ is zero. A Lipschitz composition then
has derivative zero. In $V_2$, the term differentiating the cofactor
also vanishes because $\Psi(U,0)=0$. Hence, almost everywhere on
$\{P=0\}$,
\[
 \nabla F=0,\qquad \operatorname{div}V_1
       =\operatorname{div}V_2=0.
\]
These assertions agree with the weak derivatives by Rademacher's
theorem; see \cite[\S\S3.1, 4.2]{EG}. The remaining density at
such points is $\lambda\ell^2U^2\ge0$, including when the critical
set has positive area. Integration, together with
$\int|D^2U|^2=\int|\Delta U|^2$, proves \eqref{functional}.
\end{proof}

We now prove Theorem~\ref{main}.
\begin{proof}
Let $0\ne v\in C_c^\infty(\Omega)$ and set
\[
 Q=\frac{\|\Delta v\|_2^2}{\|\nabla v\|_2^2},\qquad
 \widetilde\Omega=\sqrt Q\,\Omega,\qquad U(x)=v(x/\sqrt Q).
\]
In dimension two, $\|\Delta U\|_2^2=\|\nabla U\|_2^2$.
The function $F$ in Proposition~\ref{prop:functional} is Lipschitz
and has compact support in $\widetilde\Omega$, so it belongs to
$H_0^1(\widetilde\Omega)$ by mollification. At a nonzero maximum
point of $|U|$, one has $\nabla U=0$ and $F=\ell|U|>0$; therefore
$F$ is not identically zero. Applying \eqref{functional} gives
\[
 Q^{-1}\lambda_1^{\mathrm D}(\Omega)
   =\lambda_1^{\mathrm D}(\widetilde\Omega)
   \le\lambda=\frac{j_{0,1}^2}{j_{1,1}^2}.
\]
Taking the infimum over $v$ and then using density in $H_0^2(\Omega)$
proves \eqref{eq:comparison}. The disk eigenvalues show that the
constant is optimal. The proof only uses compactly supported tests,
so boundedness of $\Omega$ is not required.
\end{proof}

\section{Equality}\label{sec:equality}
We first isolate the distinction between equality almost everywhere
and equality of open sets.
\begin{lemma}\label{lem:strict-disk}
For every nonempty proper open subset $W$ of a planar disk $B$,
\[
 \Lambda_1(W)>\Lambda_1(B).
\]
\end{lemma}
\begin{proof}
Choose $x_0\in B\setminus W$. By Lemma~\ref{app:formfacts}, the
buckling minimum on $W$ is attained, and every element of $H_0^2(W)$,
extended by zero, has a continuous representative vanishing at $x_0$.
If the two eigenvalues were equal, a minimizer on $W$ would also
minimize on $B$. Lemma~\ref{app:disk} identifies the first eigenspace
on $B$ as one-dimensional, spanned by a function strictly positive
inside $B$. No nonzero element of this eigenspace vanishes at $x_0$.
\end{proof}

We now prove Corollary~\ref{cor:area}.
\begin{proof}
The Faber--Krahn inequality \cite{Faber,Krahn} gives
\[
 |\Omega|\lambda_1^{\mathrm D}(\Omega)\ge\pi q^2
\]
for every open set of finite positive area. Multiplication by $k^2/q^2$
and Theorem~\ref{main} give \eqref{eq:area}.

If equality holds in \eqref{eq:area}, then equality holds in the
Faber--Krahn inequality. Its equality characterization gives a disk $B$
of area $|\Omega|$ such that $\Omega$ and $B$ agree almost everywhere.
Openness now implies $\Omega\subset B$, since an open
neighborhood of any point outside $B$, including a point on its
boundary, has positive area outside $B$.
The equality of areas gives $\Lambda_1(\Omega)=\Lambda_1(B)$.
Lemma~\ref{lem:strict-disk} excludes a proper inclusion. Conversely,
a disk gives equality by its explicit first eigenfunction.
\end{proof}

\begin{remark}
For a smooth bounded simply connected planar domain, the first
Dirichlet Stokes eigenvalue equals $\Lambda_1(\Omega)$ under the
stream-function representation; see \cite[Remark~1]{HMP}.
Here the Stokes eigenvalue is the infimum of
$\int_\Omega|\nabla v|^2/\int_\Omega|v|^2$ over nonzero
$v\in H_0^1(\Omega;\mathbb R^2)$ with $\operatorname{div}v=0$.
Thus the area inequality applies to that
problem in the simply connected class. On a multiply connected
domain, a stream function can have different constants on different
boundary components and need not belong to $H_0^2(\Omega)$.
The present argument does not remove that distinction, nor does it
prove the buckling conjecture in dimensions $N\ge3$.
\end{remark}
\appendix
\section{Bessel constants and root identification}\label{app:constants}
The estimates in this appendix are finite comparisons of explicitly
stated rational polynomials and their alternating remainders.

\begin{lemma}
\label{subm:J1-separation}
If $0<a<b$ are distinct positive zeros of $J_1$, then $b-a>\pi$.
\end{lemma}
\begin{proof}
The function $y(r)=\sqrt r\,J_1(r)$ satisfies
\[
 y''+\left(1-\frac{3}{4r^2}\right)y=0\qquad(r>0).
\]
It is a nonzero member of $H_0^1(a,b)$.  Multiplication by $y$ and
integration by parts, followed by the one-dimensional Dirichlet
Wirtinger inequality, gives
\[
 \frac{\pi^2}{(b-a)^2}\int_a^b y^2\,dr
 \le\int_a^b(y')^2\,dr
 =\int_a^b\left(1-\frac{3}{4r^2}\right)y^2\,dr
 <\int_a^b y^2\,dr.
\]
Since the last integral is positive, $b-a>\pi$.
\end{proof}

\begin{lemma}

Writing $q_n=j_{0,n}$ and $\kappa_n=j_{1,n}$, one has
\begin{equation}\label{subm:root-boxes}
\begin{array}{c|c|c}
 \text{zero}&\text{rational interval}&\text{endpoint signs}\\\hline
 q=q_1&(2.40482,2.40483)&(+,-)\\
 k=\kappa_1&(3.83170,3.83171)&(+,-)\\
 q_2&(5.5200,5.5202)&(-,+)\\
 \kappa_2&(7.015,7.016)&(-,+)\\
 \kappa_3&(10.173,10.174)&(+,-).
\end{array}
\end{equation}
Each displayed interval contains exactly its indicated zero.
\end{lemma}
\begin{proof}
Here and below, the signs for $q_n$ concern $J_0$ and those for
$\kappa_n$ concern $J_1$.  For $\nu=0,1$ put
\[
 S_{\nu,N}(x)=\sum_{m=0}^N
 \frac{(-1)^m(x/2)^{2m+\nu}}{m!(m+\nu)!}.
\]
For $0<x<11$, the absolute terms in the tail starting at $m=31$
strictly decrease: their successive ratio is
\[
 \frac{x^2}{4(m+1)(m+\nu+1)}
 <\frac{121}{4\cdot32^2}<1\qquad(m\ge31).
\]
Consequently
\begin{equation}
 S_{\nu,31}(x)<J_\nu(x)<S_{\nu,30}(x),\qquad
 0<S_{\nu,30}(x)-J_\nu(x)
 <\frac{(x/2)^{62+\nu}}{31!(31+\nu)!}.
\end{equation}
The following table gives integers
$A,B$ for which
\[
 \frac A{10^8}<S_{\nu,31}(x)<S_{\nu,30}(x)<\frac B{10^8}.
\]
Every entry is a rational comparison obtained by substitution in the
two displayed finite sums and multiplication by positive denominators.
\begin{equation}\label{subm:rational-record}
\begin{array}{c|r|r|r}
 \nu&x&A&B\\\hline
 0&2.40482&288&289\\
 0&2.40483&-231&-230\\
 1&3.83170&240&241\\
 1&3.83171&-163&-162\\
 0&5.5200&-2658&-2657\\
 0&5.5202&4147&4148\\
 1&7.015&-17608&-17607\\
 1&7.016&12404&12405\\
 1&10.173&11689&11690\\
 1&10.174&-13281&-13280\\\hline
 1&2.40482&51914869&51914870\\
 0&3.83170&-40275940&-40275939\\
 1&5.5201&-34026346&-34026345\\
 0&10.1735&-24970488&-24970487
\end{array}
\end{equation}
The first ten rows give all endpoint signs in
\eqref{subm:root-boxes}, hence the existence of a zero in every
interval.

We next identify the root indices.  The alternating series gives
\[
 J_1(r)>\frac r2-\frac{r^3}{16}\ge0
 \qquad(0<r\le\sqrt8).
\]
The first inequality is strict even at $r=\sqrt8$, since the
remaining alternating tail begins with a positive term and strictly
decreasing absolute terms.  Let $a_1,a_2,a_3$ be zeros in the three
displayed $J_1$ intervals, in increasing order.  Each interval has
width less than $1<3<\pi$, so Lemma~\ref{subm:J1-separation}
makes its zero unique.  There is no positive zero before $a_1$:
such a zero would lie above $\sqrt8>2$, and its distance to $a_1$
would be less than $3.83171-2<3<\pi$.

If a further zero lay between $a_1$ and $a_2$, applying
Lemma~\ref{subm:J1-separation} to the two subintervals would give
$a_2-a_1>2\pi$.  This contradicts
\[
 a_2-a_1<7.016-3.83170=3.18430<6<2\pi.
\]
The same argument between $a_2$ and $a_3$ uses
\[
 a_3-a_2<10.174-7.015=3.159<6<2\pi.
\]
Thus $a_i=\kappa_i$ for $i=1,2,3$.  These are the only numerical
uses of $\pi$: the elementary lower bound $\pi>3$ suffices; no
upper or finer lower bound is used.

For $\nu\ge0$ the classical interlacing theorem is
\[
 j_{\nu,1}<j_{\nu+1,1}<j_{\nu,2}<j_{\nu+1,2}<\cdots
 \qquad\text{\cite[Eq.~10.21.2]{DLMF}}.
\]
For $\nu=0$ it supplies exactly one zero of $J_0$ in
$(0,\kappa_1)$ and exactly one in $(\kappa_1,\kappa_2)$.
The $J_0$ intervals in \eqref{subm:root-boxes} lie in these respective
gaps, so their zeros are $q_1$ and $q_2$, with uniqueness as asserted.
\end{proof}

\begin{lemma}\label{gm:coarse}
With $\alpha=q/k$, $\lambda=\alpha^2$, $c=-J_0(k)$, and
$K=\alpha^2J_1(q)^2/c^2$, one has
\[
 .6276<\alpha<.628,\qquad \frac14<\lambda<\frac25,
 \qquad .4027<c<.4028,\qquad \frac{19}5<k<4,
 \qquad K>.654.
\]
Furthermore
\[
 B(\kappa_3)>.634,\qquad B(q_2)<.640,\qquad
 J_1(q_2)^2<.116,
 \qquad k<\alpha\kappa_2<q_2<\alpha\kappa_3.
\]
\end{lemma}
\begin{proof}
The integer-order integral representation
\cite[Eq.~10.9.2]{DLMF}, differentiated for real $x$, gives
\[
 J_n(x)=\frac1\pi\int_0^\pi
       \cos(x\sin\vartheta-n\vartheta)\,d\vartheta,
 \qquad |J_n'(x)|\le1\quad(n\in\mathbb Z).
\]
In particular, the last four rows of
\eqref{subm:rational-record}, the root intervals, and the mean-value
theorem yield
\[
 J_1(q)>.51914-.00001>.5191,\qquad
 -.4028<J_0(k)<-.4027,
\]
\[
 |J_1(q_2)|<.34027+.0001<.3404,\qquad
 J_0(10.1735)<-.2497.
\]
For example, the asserted two-sided bound at $k$ follows already from
\[
 -.40275940-.00001<J_0(k)
                  <-.40275939+.00001.
\]
The root intervals give
\[
 .6276\cdot3.83171<2.40482,
 \qquad 2.40483<.628\cdot3.83170,
\]
and hence the bounds on $\alpha$.  Since $q>2$ and $k<4$, one has
$\lambda>1/4$; and $.628^2<2/5$ proves its upper bound.  Now
\[
 K>\frac{(.6276)^2(.5191)^2}{(.4028)^2}
 =\frac{66335796013041}{101404900000000}>.654,
\]
\[
 J_1(q_2)^2<(.3404)^2=.11587216<.116,
\]
\[
 B(q_2)=q_2J_1(q_2)^2
 <5.5202(.3404)^2=.639637497632<.640.
\]
On the interval $(10.173,10.174)$, uniqueness of the zero and the
endpoint signs make $J_1$ positive before $\kappa_3$ and negative
after it.  Thus $J_0'=-J_1$ makes $\kappa_3$ the minimum of $J_0$
on the entire closed interval, and therefore
\[
 J_0(\kappa_3)\le J_0(10.1735)<-.2497.
\]
It follows that
\[
 B(\kappa_3)=\kappa_3J_0(\kappa_3)^2
 >10.173(.2497)^2=.63428746557>.634.
\]
Finally the required order comparisons follow from the explicit
rational chains
\[
\begin{aligned}
 k&<3.83171<.6276\cdot7.015=4.402614<\alpha\kappa_2,\\
 \alpha\kappa_2&<.628\cdot7.016=4.406048<5.5200<q_2,\\
 q_2&<5.5202<.6276\cdot10.173=6.3845748<\alpha\kappa_3.
\end{aligned}
\]
\end{proof}

The same rational series bounds supply the few small-argument
comparisons used earlier.  At $r=19/10$,
\[
 \frac7{25}<S_{0,3}(r)<J_0(r)<S_{0,4}(r)
 <\frac{283}{1000},\qquad
 \frac{J_1(r)}r>\frac{S_{1,3}(r)}r>\frac3{10}.
\]
Also
\[
 J_0(2)>S_{0,3}(2)=\frac29>\frac{11}{50},\qquad
 J_1(2)<S_{1,2}(2)=\frac7{12}<\frac35.
\]
All needed alternating tails are decreasing from their first omitted
term.  These estimates prove the asserted coarse bounds.

\section{Zero extension, compactness and the disk spectrum}

\begin{lemma}\label{app:formfacts}
Let $\Omega\subset\mathbb R^2$ be bounded and open. Zero extension
identifies $H_0^m(\Omega)$ isometrically with a closed subspace of
$H_0^m(Q)$, for $m=1,2$, whenever $Q$ is a rectangle containing
$\overline\Omega$ in its interior. There is a constant $C_Q$ such that
\[
 \|u\|_{H^2(\Omega)}\le C_Q\|\Delta u\|_{L^2(\Omega)},
 \qquad u\in H_0^2(\Omega).
\]
The embedding $H_0^2(\Omega)\hookrightarrow H_0^1(\Omega)$ is compact.
Every element of $H_0^2(\Omega)$ has a continuous zero extension to
$\mathbb R^2$, with bounded point evaluations. The buckling infimum is attained whenever $\Omega$ is nonempty.
\end{lemma}
\begin{proof}
Zero extension preserves the $H^m$ norm on $C_c^\infty(\Omega)$;
passing to the completion proves the first assertion and closedness.
For smooth compactly supported functions, integration by parts gives
\[
 \|D^2u\|_2^2=\|\Delta u\|_2^2,\qquad
 \|\nabla u\|_2^2=-\int u\Delta u.
\]
Poincar\'e's inequality on $Q$ first bounds $\|\nabla u\|_2$ by
$C_Q\|\Delta u\|_2$, and then bounds $\|u\|_2$. Density gives the
displayed coercive estimate. Rellich compactness on the fixed
rectangle $Q$ gives strong $H^1(Q)$ convergence of a subsequence
from every bounded $H^2(Q)$ sequence, hence the asserted compactness.
These are the usual embedding and compactness theorems, applied only
on a rectangle; see \cite[Chapters~4 and~6]{AF}.

For completeness, point evaluation follows directly on the whole
plane. With any fixed Plancherel normalization of the Fourier
transform, Cauchy--Schwarz gives
\[
 \|\widehat u\|_{L^1(\mathbb R^2)}
 \le\left(\int_{\mathbb R^2}(1+|\xi|^2)^{-2}\,d\xi\right)^{1/2}
       \|(1+|\xi|^2)\widehat u\|_2
 \le C\|u\|_{H^2(\mathbb R^2)}.
\]
The first integral is finite. Fourier inversion therefore gives a
continuous representative and $\|u\|_\infty\le C\|u\|_{H^2}$.
This also proves uniform convergence of continuous representatives
under $H^2$ convergence. Approximating by $C_c^\infty(\Omega)$
therefore shows that this representative is zero at every point
of $\mathbb R^2\setminus\Omega$, including its boundary.
Finally, normalize the gradient energy of
a buckling minimizing sequence to one. The coercive estimate, weak
$H^2$ compactness and strong $H^1$ compactness give a minimizer with
unit gradient energy. No extension theorem for the boundary of
$\Omega$ is used.
\end{proof}

\begin{lemma}\label{app:disk}
On a disk $B_R(a)$, the first buckling eigenvalue is $k^2/R^2$,
where $k=j_{1,1}$. Its eigenspace is one-dimensional and spanned by
\[
 U_B(x)=J_0(k|x-a|/R)-J_0(k),
\]
which is strictly positive in the disk.
\end{lemma}
\begin{proof}
Translation reduces to $a=0$. A weak buckling eigenfunction satisfies
$\Delta(\Delta+\Lambda)U=0$. The distribution
$(\Delta+\Lambda)U$ is harmonic, hence smooth; interior elliptic
regularity for the second-order equation then makes $U$ smooth
in the disk. The relevant local regularity results can be found in
\cite[Sections~2.2 and~6.3]{EvansPDE}.

Put $\rho=\sqrt\Lambda>0$. The regular solution in angular mode
$n\ge0$ is a linear combination of $r^n$ and $J_n(\rho r)$,
with $r^0=1$ for $n=0$. Indeed applying $\Delta+\Lambda$ to its
regular harmonic component gives a particular solution proportional
to $r^n$, and the remaining regular Helmholtz component is
$J_n(\rho r)$. The two clamped traces at $r=R$ have a nonzero
solution precisely when
\[
 \rho R J_n'(\rho R)-nJ_n(\rho R)
             =-\rho R J_{n+1}(\rho R)=0.
\]
The trace assertion is the classical trace theorem on the smooth
disk, applied to $U$ and its first derivatives, which belong to
$H_0^1(B_R)$ by the $H_0^2$ closure definition; see
\cite[Section~4.3]{EG}. Each separated radial solution is
smooth up to $r=R$, so these are ordinary endpoint values.
The angular Fourier projections are bounded on $H^2(B_R)$: they
are averages of rotations, which preserve the equivalent squared
Sobolev norm $\|u\|_2^2+\|\nabla u\|_2^2+\|D^2u\|_2^2$.
They preserve $H_0^2(B_R)$ by its closure definition, commute with
the weak equation, and preserve the zero traces. Thus every
nonzero Fourier component of an eigenfunction is subject to the
displayed characteristic equation, and a nonzero eigenfunction
has at least one such component.
Interlacing of the positive zeros of consecutive Bessel orders gives
$j_{1,1}<j_{2,1}<j_{3,1}<\cdots$; see
\cite[Section~10.21(i)]{DLMF}. All nonzero zeros are simple.
Thus the smallest admissible $\rho R$ is the simple zero $k$ in
mode $n=0$, which yields exactly the displayed function.
The mode calculation also shows that no other angular component
can belong to this eigenspace. The displayed $U_B$ is an admissible
clamped test function: near the boundary it satisfies
$U_B=O((R-r)^2)$ and $\nabla U_B=O(R-r)$. Multiply it by a smooth
cutoff that is zero for $R-r\le\varepsilon$ and one for
$R-r\ge2\varepsilon$. The error in each second derivative is
uniformly bounded and supported in a collar of area $O(\varepsilon)$,
so the cutoffs converge to $U_B$ in $H^2$. Hence $U_B\in H_0^2(B_R)$.
Its differential equation and integration by parts give its
buckling quotient $k^2/R^2$, proving attainment of the stated
smallest value. Finally $J_1>0$ on $(0,k)$ and
$J_0'=-J_1$ imply its strict interior positivity.
\end{proof}

\begin{thebibliography}{99}

\bibitem{AF}
R.~A. Adams and J.~J.~F. Fournier,
\emph{Sobolev Spaces}, 2nd ed.,
Pure and Applied Mathematics, vol.~140,
Academic Press, Amsterdam, 2003.

\bibitem{AB}
M.~S. Ashbaugh and D. Bucur,
On the isoperimetric inequality for the buckling of a clamped plate,
\emph{Z. Angew. Math. Phys.} \textbf{54} (2003), no.~5, 756--770.
\doi{10.1007/s00033-003-3204-3}.

\bibitem{AL}
M.~S. Ashbaugh and R.~S. Laugesen,
Fundamental tones and buckling loads of clamped plates,
\emph{Ann. Scuola Norm. Sup. Pisa Cl. Sci. (4)}
\textbf{23} (1996), no.~2, 383--402.

\bibitem{BramblePayne}
J.~H. Bramble and L.~E. Payne,
Pointwise bounds in the first biharmonic boundary value problem,
\emph{J. Math. Phys.} \textbf{42} (1963), 278--286.
\doi{10.1002/sapm1963421278}.

\bibitem{BP}
D. Buoso and E. Parini,
The buckling eigenvalue problem in the annulus,
\emph{Commun. Contemp. Math.} \textbf{23} (2021), no.~4, 2050044.
\doi{10.1142/S0219199720500443}.

\bibitem{EvansPDE}
L.~C. Evans,
\emph{Partial Differential Equations}, 2nd ed.,
Graduate Studies in Mathematics, vol.~19,
American Mathematical Society, Providence, RI, 2010.

\bibitem{EG}
L.~C. Evans and R.~F. Gariepy,
\emph{Measure Theory and Fine Properties of Functions}, revised ed.,
Textbooks in Mathematics, CRC Press, Boca Raton, FL, 2015.

\bibitem{Faber}
G. Faber,
Beweis, da\ss{} unter allen homogenen Membranen von gleicher Fl\"ache
und gleicher Spannung die kreisf\"ormige den tiefsten Grundton gibt,
\emph{Sitzungsber. Bayer. Akad. Wiss. Math.-Phys. Kl.}
(1923), 169--172.

\bibitem{Friedlander}
L. Friedlander,
Remarks on the membrane and buckling eigenvalues for planar domains,
\emph{Mosc. Math. J.} \textbf{4} (2004), no.~2, 369--375.
\doi{10.17323/1609-4514-2004-4-2-369-375}.

\bibitem{HMP}
A. Henrot, I. Mazari-Fouquer and Y. Privat,
Is the Faber--Krahn inequality true for the Stokes operator?,
\emph{Calc. Var. Partial Differential Equations}
\textbf{63} (2024), article~228.
\doi{10.1007/s00526-024-02820-7}.

\bibitem{Krahn}
E. Krahn,
\"Uber eine von Rayleigh formulierte Minimaleigenschaft des Kreises,
\emph{Math. Ann.} \textbf{94} (1925), 97--100.
\doi{10.1007/BF01208645}.

\bibitem{LN}
J. Lamboley and M. Nahon,
Boundary regularity of a fourth order Alt-Caffarelli problem and
applications to the minimization of the critical buckling load,
preprint, 2025, arXiv:2512.18626.
\doi{10.48550/arXiv.2512.18626}.

\bibitem{LP}
M. Levitin and I. Polterovich,
Payne's conjecture for buckling eigenvalues of odd index,
preprint, 2026, arXiv:2609.12235.
\doi{10.48550/arXiv.2609.12235}.

\bibitem{DLMF}
National Institute of Standards and Technology,
\emph{NIST Digital Library of Mathematical Functions},
version~1.2.8, released 15 September 2026,
\url{https://dlmf.nist.gov/}, accessed 15 October 2026.

\bibitem{Payne}
L.~E. Payne,
Inequalities for eigenvalues of membranes and plates,
\emph{J. Rational Mech. Anal.} \textbf{4} (1955), 517--529.

\bibitem{PS}
G. P\'olya and G. Szeg\H{o},
\emph{Isoperimetric Inequalities in Mathematical Physics},
Annals of Mathematics Studies, vol.~27,
Princeton University Press, Princeton, NJ, 1951.

\bibitem{Stollenwerk}
K. Stollenwerk,
On the optimal domain for minimizing the buckling load of a clamped plate,
\emph{Z. Angew. Math. Phys.} \textbf{74} (2023), article~12.
\doi{10.1007/s00033-022-01898-3}.

\bibitem{SW}
K. Stollenwerk and A. Wagner,
Optimality conditions for the buckling of a clamped plate,
\emph{J. Math. Anal. Appl.} \textbf{432} (2015), no.~1, 254--273.
\doi{10.1016/j.jmaa.2015.06.035}.

\bibitem{Szego50}
G. Szeg\H{o},
On membranes and plates,
\emph{Proc. Natl. Acad. Sci. USA} \textbf{36} (1950), no.~3, 210--216.
\doi{10.1073/pnas.36.3.210}.

\bibitem{Szego58}
G. Szeg\H{o},
Note to my paper ``On membranes and plates'',
\emph{Proc. Natl. Acad. Sci. USA} \textbf{44} (1958), no.~4, 314--316.
\doi{10.1073/pnas.44.4.314}.

\end{thebibliography}
\end{document}
